% !TeX program = XeLaTeX
% !TeX root = gitabook.tex
\documentclass[oneside,12pt]{book}
\usepackage{etoolbox}

\newbool{kindle}
\newbool{print}
\setbool{kindle}{true}
\setbool{print}{false}
\ifbool{kindle}{\usepackage[paperwidth=192mm,paperheight=254mm,left=4mm,right=4mm,top=9mm,bottom=10mm]{geometry}\addtolength{\headsep}{-4mm}}{\usepackage[top=1.5cm, bottom=1.8cm, left=1.5cm, right=1.5cm,paperwidth=148mm,paperheight=210mm]{geometry}}
\usepackage{hypershloka}
\input{../preamble}
\ifbool{print}{%Margin changes for print
\addtolength{\evensidemargin}{-0.5cm}
\addtolength{\oddsidemargin}{0.5cm}}{}

\begin{document}
\setmainfont[Script=Devanagari]{Siddhanta}
\input{../frontmatter}
\input{../preface.tex}


\newcommand{\ganapatyadiDhyanam}{

\dnsub{श्री-महागणपति-प्रार्थना}

\twolineshloka*
{शुक्लाम्बरधरं विष्णुं शशिवर्णं चतुर्भुजम्}
{प्रसन्नवदनं ध्यायेत् सर्वविघ्नोपशान्तये}

\twolineshloka*
{वागीशाद्याः सुमनसः सर्वार्थानामुपक्रमे}
{यं नत्वा कृतकृत्याः स्युस्तं नमामि गजाननम्}

\dnsub{श्री-गुरु-प्रार्थना}

\twolineshloka*
{गुरुर्ब्रह्मा गुरुर्विष्णुर्गुरुर्देवो महेश्वरः}
{गुरुः साक्षात् परं ब्रह्म तस्मै श्री-गुरवे नमः}

\twolineshloka*
{सदाशिवसमारम्भां शङ्कराचार्यमध्यमाम्}
{अस्मदाचार्यपर्यन्तां वन्दे गुरुपरम्पराम्}

\dnsub{श्री-सरस्वती-प्रार्थना}
\fourlineindentedshloka*
{दोर्भिर्युक्ता चतुर्भिः स्फटिकमणिनिभैरक्षमालां दधाना}
{हस्तेनैकेन पद्मं सितमपि च शुकं पुस्तकं चापरेण}
{भासा कुन्देन्दुशङ्खस्फटिकमणिनिभा भासमानाऽसमाना}
{सा मे वाग्देवतेयं निवसतु वदने सर्वदा सुप्रसन्ना}

\dnsub{श्री-व्यास-नमस्क्रिया}
\resetShloka

\twolineshloka
{व्यासं वसिष्ठनप्तारं शक्तेः पौत्रमकल्मषम्}
{पराशरात्मजं वन्दे शुकतातं तपोनिधिम्}

\fourlineindentedshloka
{अभ्रश्यामः पिङ्गजटाबद्धकलापः}
{प्रांशुर्दण्डी कृष्णमृगत्वक्परिधानः}
{साक्षाल्लोकान् पावयमानः कविमुख्यः}
{पाराशर्यः पर्वसु रूपं विवृणोतु}

\twolineshloka
{जयति पराशरसूनुः सत्यवतीहृदयनन्दनो व्यासः}
{यस्यास्यकमलगलितं वाङ्मयममृतं जगत् पिबति}

\resetShloka
}

\newcommand{\genMangalaShloka}{
\dnsub{मङ्गलश्लोकाः}
\fourlineindentedshloka*
{स्वस्ति प्रजाभ्यः परिपालयन्तां}
{न्यायेन मार्गेण महीं महीशाः}
{गोब्राह्मणेभ्यः शुभमस्तु नित्यं}
{लोकाः समस्ताः सुखिनो भवन्तु}

\twolineshloka*
{काले वर्षतु पर्जन्यः पृथिवी सस्यशालिनी}
{देशोऽयं क्षोभरहितो ब्राह्मणाः सन्तु निर्भयाः}

\twolineshloka*
{अपुत्राः पुत्रिणः सन्तु पुत्रिणः सन्तु पौत्रिणः}
{अधनाः सधनाः सन्तु जीवन्तु शरदां शतम्}

\twolineshloka*
{यदक्षरपदभ्रष्टं मात्राहीनं तु यद् भवेत्}
{तत् सर्वं क्षम्यतां देव नारायण नमोऽस्तुते}

\twolineshloka*
{विसर्गबिन्दुमात्राश्च पदपादाक्षराणि च}
{न्यूनानि चातिरिक्तानि क्षमस्व पुरुषोत्तम}

\twolineshloka*
{यज्ञेशाच्युत गोविन्द माधवानन्त केशव}
{कृष्ण विष्णो हृषीकेश वासुदेव नमोऽस्तुते}

\fourlineindentedshloka*
{कायेन वाचा मनसेन्द्रियैर्वा}
{बुद्‌ध्याऽऽत्मना वा प्रकृतेः स्वभावात्}
{करोमि यद्यत् सकलं परस्मै}
{नारायणायेति समर्पयामि}}

\newcommand{\padmaPuranam}{
\dnsub{मङ्गलश्लोकः}
\fourlineindentedshloka*
{स्वच्छं चन्द्रावदातं करिकर-मकर-क्षोभ-सञ्जात-फेनं}
{ब्रह्मोद्भूतिप्रसक्तैर्व्रतनियमपरैः सेवितं विप्रमुख्यैः}
{ओङ्कारालङ्कृतेन त्रिभुवनगुरुणा ब्रह्मणा दृष्टिभूतं}
{सम्भोगा भोग रम्यं जलमशुभहरं पौष्करो नः पुनातु }
}


\part{मूलम्}
\padacchedafalse  % Ensure we are in Moola mode
% \setmainfont[Script=Devanagari,FakeStretch=0.9]{Siddhanta}
% \setmainfont[Script=Devanagari]{Sanskrit 2003}
\setcounter{page}{0}
\pagenumbering{arabic}
\sectionmark{\mbox{}}
\begin{center}
\fontsize{19pt}{23pt}\selectfont
\input{../nyasa.tex}
\clearpage
\phantomsection\label{sec:start_moola}
\input{../gita.tex}
\closesection
\input{../mahatmyam.tex}
\input{../mahatmyam-varaha-puranam.tex}
\nobreak\hyperref[sec:start_moola]{\closesection}
\end{center}


\part{पदच्छेदः}
\padacchedatrue  % Ensure we are in Moola mode
% \input{frontmatter}
% \setmainfont[Script=Devanagari,FakeStretch=0.9]{Siddhanta}
% \setmainfont[Script=Devanagari]{Sanskrit 2003}
% \setcounter{page}{0}
% \pagenumbering{arabic}
% \sectionmark{\mbox{}}
\begin{center}
\fontsize{19pt}{23pt}\selectfont
\stepcounter{chapter}
\input{../nyasa.tex}
\clearpage
\phantomsection\label{sec:start_split}
\input{gita-words.tex}
\closesection
\input{../mahatmyam.tex}
\input{../mahatmyam-varaha-puranam.tex}
\nobreak\hyperref[sec:start_split]{\closesection}
\end{center}

\newpage\mbox{}
\clearemptydoublepage

\part{पद्मपुराणान्तर्गत-गीता-माहात्म्यम्}
\begin{center}
\fontsize{19pt}{23pt}\selectfont
\input{../mahatmyam-padma-puranam.tex}
\end{center}


\part{श्लोकानुक्रमणिका}
\begingroup
\setmainfont[Script=Devanagari,Mapping=devanagarinumerals,Scale=1.2]{Adishila}
\fontsize{16pt}{18pt}\selectfont
\input{index_moola.tex}

\closesub

\endgroup

\part{पदानुक्रमणिका}
\begingroup
\setmainfont[Script=Devanagari,Mapping=devanagarinumerals,Scale=1.2]{Adishila}
\fontsize{16pt}{18pt}\selectfont
\input{index_word.tex}

\closesub

\endgroup

\end{document}